\section{The Tool-Lab Methodology}

Tool-Lab adapts process tracing to subjects whose only interface to the world is
linguistic and tool-mediated. Instead of watching a human participant move a
cursor over a visual decision board, we observe an LLM interacting through a
controlled action space: it can deliberate internally, call tools to inspect
currently available information, and emit text outputs such as intermediate
judgments or final choices. The experimental trace is therefore defined over
tool-mediated acquisitions and decisions rather than over mouse movements.
This preserves the core Mouselab idea that the order and extent of information
acquisition are themselves behavioral data.

Crucially, dynamic tasks need not be reduced to Payne-style static matrices in
order to be implemented for LLMs. The environment can maintain a latent state
and expose only the subset of cues that is currently available, updating that
state after each action. In this way, Tool-Lab can preserve transient cue
availability, missed-information costs, staged decisions, and other features of
dynamic tasks such as election campaigns, while still logging a clean and fully
structured action trace.

Resource scarcity can also be implemented directly in monetary terms. At the
start of an episode, the system prompt specifies a maximum budget $B$ (e.g.,
\$0.10). As the model reasons and invokes tools, the environment debits the
actual billed cost of model inference and tool execution under the prevailing
pricing schedule. This cost signal can be computed from observed usage without
recording the semantic contents of private reasoning. Each tool response then
returns not only the requested information, but also the marginal cost of that
action, cumulative expenditure, and the budget remaining. Monetary budget thus
serves as a domain-general implementation of search cost for LLM subjects,
playing the role that time pressure, effort, or access difficulty plays in
human process-tracing paradigms.